% Section 1: the measurement problem. This section exists because the honest
% answer to "what do Austin musicians earn" depends on which instrument you ask,
% and the instruments disagree in a patterned way. It opens with an orientation
% subsection that nests Austin inside Texas and Texas inside the country, so
% that the Austin figures in later sections have a reference point. That
% orientation only points to the comparisons; each one is developed in the
% section named in the table.
\section{Austin stopped asking musicians what they earn}
\label{sec:measurement}

What Austin musicians earn depends on which source you ask, and the sources
disagree in a patterned way. This section opens by placing Austin inside Texas
and Texas inside the country, then traces what the city measured and then
stopped measuring, follows one widely repeated income figure back to its origin,
sets out what four federal statistical programs can and cannot see, and ends
with the rules this study follows when sources conflict.

% Reserve enough space for the heading, the orientation paragraph and the
% nesting table together. The paragraph promises a table the reader should be
% able to see without turning the page, and holding the block together costs no
% page here: the study runs to the same length either way.
\needspace{30\baselineskip}
\subsection{Austin is a concentrated music town inside a state that commits
little to music}

A reader meeting an Austin figure in this study needs two reference points:
where Austin stands within Texas, and where Texas stands within the country.
Austin accounts for \NPumsAustinShareOfTxMusicians\ of Texas musicians and
singers against \NPumsAustinShareOfTxWorkers\ of employed Texans,\footnote{The
Census Bureau's American Community Survey Public Use Microdata Sample is the
only federal instrument that identifies a respondent's detailed occupation and
attaches survey weights that project the sample to a state or metro population.
The Bureau makes the raw records available for reanalysis, and both shares
here are ratios of person weights from the 2020--2024 five-year Texas file at
\url{https://www2.census.gov/programs-surveys/acs/data/pums/2024/5-Year/}.
Population: Texas residents 18--64 who worked in the reference week (ESR codes
1, 2, 4 and 5) and reported positive earnings; the musician cell restricts
that population to Musicians and Singers, occupation code OCCP 2752. Both
shares are computed on the same weighted denominator and so cancel any common
weighting artefact.} so the metro's share of the state's musicians is close to
twice its share of the state's workers. Travis County registers more solo
creative businesses per resident than Dallas, Bexar or Harris, the three other
large Texas counties compared here.\footnote{The Census Bureau's Nonemployer
Statistics program tabulates every U.S.\ business with no paid employees that
files a federal income-tax return with receipts of at least \$1,000, and
publishes county-by-industry counts at
\url{https://www.census.gov/programs-surveys/nonemp.html}. The count used here
is NAICS 7115, Independent Artists, Writers and Performers, per 10,000
residents for 2023, with resident totals from the Census Bureau's Vintage 2025
Population Estimates Program. That industry code covers writers and other
independent performers as well as musicians, so the count is a proxy for solo
creative work rather than a musician-only census; the same NAICS boundary
applies to every county compared, so the ranking is internally consistent.}
Texas itself has a small arts economy for its size and pays out little per
resident through its main music-specific incentive, though in the one rebate
round with a published total the Texas Music Office awarded the full amount
the state dedicates to the program each year.\footnote{The Bureau of Economic Analysis measures the arts-economy
share of state output through its Arts and Cultural Production Satellite
Account (ACPSA), which reweights the national GDP accounts to isolate arts and
culture value added; program page at
\url{https://www.bea.gov/data/special-topics/arts-and-culture}. The Texas
music-incentive figure is the Texas Music Incubator Rebate Program awarded in
FY2026, per the Texas Music Office program page at
\url{https://gov.texas.gov/music/tmir}. Both figures appear later with method
and vintage in Section~\ref{sec:state}.} The table below sets each comparison
beside the section that develops it, where the margins, methods and caveats
are given in full.

\vspace{3pt}
{\footnotesize
\setlength{\tabcolsep}{4pt}
\renewcommand{\arraystretch}{1.05}
\noindent\begin{tabular}{@{}P{0.28\textwidth} R{0.14\textwidth} P{0.54\textwidth}@{}}
\toprule
\textbf{Measure} & \textbf{Austin or Texas} & \textbf{Comparison, and the section that develops it} \\
\midrule
\multicolumn{3}{@{}l@{}}{\sffamily\bfseries\color{amnavy}Austin within Texas} \\
Share of Texas musicians and singers who work in the Austin metro
& \NPumsAustinShareOfTxMusicians
& The metro's share of all employed Texans is
\NPumsAustinShareOfTxWorkers\ \mbox{(Section~\ref{sec:earnings})} \\
\rowrule
Independent-artist businesses per 10,000 residents, Travis County, 2023
& \NNesEstabPerOneZerokTravisTwoZeroTwoThree
& Dallas \NNesEstabPerOneZerokDallasTwoZeroTwoThree, Bexar
\NNesEstabPerOneZerokBexarTwoZeroTwoThree, Harris
\NNesEstabPerOneZerokHarrisTwoZeroTwoThree\ \mbox{(Section~\ref{sec:payroll})} \\
\rowrule
Median earnings in the pooled music occupations (musicians, singers, music
directors and composers), Austin metro
& \NPumsMetroAustinMusicMedpernp
& Dallas-Fort Worth \NPumsMetroDfwMusicMedpernp, Houston
\NPumsMetroHoustonMusicMedpernp; Austin has the lowest point estimate of the
three and the margins overlap, so no gap is statistically firm
\mbox{(Section~\ref{sec:earnings})} \\
\midrule
\multicolumn{3}{@{}l@{}}{\sffamily\bfseries\color{amnavy}Texas within the nation} \\
Arts and cultural production as a share of economic output, 2023
& \NBeaGdpShareTxTwoZeroTwoThree
& \NBeaGdpShareUsAvgTwoZeroTwoThree\ nationally; Texas ranks
\NBeaGdpShareTxNationalRank, and of the eight benchmark states only Louisiana
(\NBeaGdpShareLaTwoZeroTwoThree) is lower \mbox{(Section~\ref{sec:state})} \\
\rowrule
State music money awarded, claimed or appropriated, per resident
& \NMusicDisbursedPercapitaTexas
& New York \NMusicDisbursedPercapitaNewyork, Tennessee
\NMusicDisbursedPercapitaTennessee, Louisiana
\NMusicDisbursedPercapitaLouisiana, Georgia
\NMusicDisbursedPercapitaGeorgia\ \mbox{(Section~\ref{sec:state})} \\
\rowrule
Federal shuttered-venue grants, Austin area
& \NSvogAustinTotal
& Dallas-Fort Worth \NSvogDfwAreaTotal, Houston \NSvogHoustonAreaTotal,
Nashville \NSvogNashvilleTotal\ \mbox{(Section~\ref{sec:state})} \\
\bottomrule
\end{tabular}
}

\srcnote{Sources: American Community Survey 2020--2024 five-year public use
microdata, employed adults 18 to 64 with positive earnings, for the shares and
the medians; Census Bureau Nonemployer Statistics, NAICS 7115, with Census
Bureau population estimates, for businesses per 10,000 residents; Bureau of
Economic Analysis Arts and Cultural Production Satellite Account for the output
share; state statutes and agency reports, with Census Bureau population
estimates, for money disbursed per resident; Small Business Administration
award file, July 2022 vintage, for the shuttered-venue grants. How to read it:
the middle column is the Austin or Texas figure and the right column is what
that figure is measured against, with the section that develops the comparison
in parentheses. Caveats: the three metro medians rest on small samples and
carry wide margins of error; independent-artist businesses cover writers and other performers as well
as musicians, so that count is a proxy for solo creative work rather than a
count of musicians; the shuttered-venue metro groupings are built from
recipient city names rather than official metropolitan boundaries, and they are
not built alike, since the two Texas figures use multi-place lists while the
Nashville figure covers the consolidated city only and so understates its metro;
\textbf{the Austin and Houston totals should not be read as a ranking}, because
several cinema-chain awards filed at Austin-area corporate addresses fund
operations in other states, and removing the two clearest of them alone puts
Austin below Houston (Section~\ref{sec:state});
shuttered-venue awards are nominal dollars, an exception to the deflation rule
below, because every award falls in one program period; the per-resident state
figures are money that left the treasury or was awarded, claimed or
appropriated, rather than dollars authorised, which is why Georgia appears at
zero; and the Tennessee figure combines film, television and music scoring
because no music-only breakout exists, so it is an upper bound rather than a
music-specific number.}

\vspace{4pt}

Those are the two reference points. The rest of this section asks how much
weight any one of those Austin figures can carry.

\subsection{The 2022 music census dropped the income questions the 2014 one
asked}

The City of Austin's Music and Entertainment Division commissioned Titan Music
Group to field a music census in 2014, and that instrument asked what people
earned from music. It drew 3,968 self-selected
responses.\footnote{Titan Music Group, LLC (for the City of Austin Economic
Development Department, Music and Entertainment Division), \textit{The Austin
Music Census: A Data-Driven Assessment of Austin's Commercial Music Economy},
fielded late 2014, published 1 June 2015, with income questions covering
calendar year 2013;
\url{https://austin.widen.net/content/om7zqn1qbx/pdf/Austin_Music_Census_Interactive_PDF_53115.pdf}.
Recruitment ran through Austin music-industry mailing lists and social channels,
so respondents self-selected into the sample rather than being drawn by
probability method; the census itself flags that limitation, and the
microdata were never publicly released. Every share below is a sample share,
not a population estimate.} Of the 1,883 respondents who identified as
musicians, 68.4\%, or 1,288 people, reported earning under \$10,000 from
music, a threshold worth \NPumsMusiccensusTenkTwoZeroTwoFiveusd\ in 2025
dollars.\footnote{The 68.4~percent figure appears on page 74 of the Titan
report; it is the share of the 1,883 self-identified musician respondents who
placed themselves in the ``under \$10,000'' band of the pre-tax
calendar-2013 earnings question. Converted with the Bureau of Labor
Statistics's Consumer Price Index for All Urban Consumers, U.S.\ city average,
all items, series CUUR0000SA0 (annual-average factor
\NPumsDeflTwoZeroOneThreeToTwoZeroTwoFive\ between 2013 and 2025), retrieved
from the BLS Public Data API at
\url{https://api.bls.gov/publicAPI/v2/timeseries/data/}. Inflation adjustment
matters more here than usual: a reader comparing the 2013 threshold against
present-day earnings without adjusting would understate today's bar by about
28 percent, since \$10,000 is that much below the
\NPumsMusiccensusTenkTwoZeroTwoFiveusd\ the line is worth now.} The 2022
successor census dropped those income questions. Its authors explain why in
the front matter: working ``with the guidance of the community,'' they removed
questions from the 2014 survey ``that aimed to quantify the income of music
people in a way that turned off many
respondents.''\footnote{Sound Music Cities, \textit{Greater Austin Music
Census, Summary Report}, published February 2023, printed page 2;
\url{https://www.austinmusiccensus.org/s/2022-Greater-Austin-Music-Census-SUMMARY-REPORT-v13-optimized.pdf}.
The team removed the income-quantification items specifically, not every
income-related question, and the 2022 questionnaire kept items on the
composition of music income, on spending, on housing and on whether
respondents expected to remain in Austin, all of which Section~\ref{sec:census}
uses. Like the 2014 fielding, the 2022 census is a self-selected online
convenience sample.}

The census's own cohort shows what the same instrument produces where an
income question stayed in. Sound Music Cities fields versions of this
questionnaire across a national cohort of music cities, and the member cities
that kept an income item publish figures: Nashville's census reports
\NCitycensusIncomeNashville\ and Baltimore's \NCitycensusIncomeBaltimore,
with the published figures running down to \NCitycensusIncomeAnchorage\ in
Anchorage.\footnote{As published by each city's census report: Nashville
\NCitycensusIncomeNashville, Baltimore \NCitycensusIncomeBaltimore, Charlotte
\NCitycensusIncomeCharlotte, Minneapolis \NCitycensusIncomeMinneapolis,
Northwest Arkansas \NCitycensusIncomeNwarkansas, Baton Rouge
\NCitycensusIncomeBatonRouge, New Orleans \NCitycensusIncomeNewOrleans,
Washington D.C.\ \NCitycensusIncomeWashingtonDc, Cleveland
\NCitycensusIncomeCleveland, Detroit \NCitycensusIncomeDetroit, Columbus
\NCitycensusIncomeColumbus, Greensboro \NCitycensusIncomeGreensboro, Tulsa
\NCitycensusIncomeTulsa, Lawrence \NCitycensusIncomeLawrence, and Anchorage
\NCitycensusIncomeAnchorage. Each figure is that city's published headline
for annual music-related income, computed on a self-selected local sample in
that city's fielding year; most of these are averages rather than medians,
and averages run above medians on right-skewed income, so the list shows what
cohort cities publish rather than a like-for-like ranking. Citations for the
city reports are in the references.} Austin has published
\NCitycensusAustinNoData. The cohort therefore shows that a workable income
item exists and that peer cities publish its results; what lapsed in Austin
is the asking, not the instrument.

The 2022 census's silence on earnings is narrower than it first appears,
because measurement of Austin musician pay continued through other
channels. The City of Austin bought a commercial series that put median
musician earnings at \$17.49 an hour in 2016 and \$17.68 in
2017.\footnote{City of Austin open data, \textit{Median Earnings of Creative
Sector Occupations}, dataset \texttt{2qxc-8cme}, at
\url{https://data.austintexas.gov/resource/2qxc-8cme.csv}, retrieved 1 August
2026. The city purchased the underlying figures from the Creative Vitality
Suite, a commercial subscription database maintained by the labor-market data
vendor Lightcast (formerly Emsi Burning Glass) that combines federal
Occupational Employment and Wage Statistics with self-employment tax records.
The vendor publishes no methodology detail on how the two sources are joined,
so a reader cannot audit the resulting hourly figure. The city's series
covers Musicians and Singers for 2016 and 2017 only.} A University of Texas
at Austin research project interviewed roughly sixty Austin
musicians and reported that their annual incomes clustered around \$30,000 to
\$35,000.\footnote{University of Texas at Austin,
\textit{Support Austin Musicians}, project website drafted July 2024,
\url{https://supportaustinmusicians.org/final-summary/}. Recruitment ran
through Austin music-community contacts rather than by any probability
method, and the public write-up summarises about sixty interviews. This
study takes the direction of the finding rather than the point estimate from
that source.} The Health Alliance for Austin Musicians,
a local nonprofit that provides health-care access to working Austin
musicians, publishes an annual impact report that names the household-income
ceiling its clients must fall under to
qualify.\footnote{Health Alliance for Austin Musicians (HAAM),
\textit{Annual Impact Report} series, 2014 through 2025, at
\url{https://www.myhaam.org/}. HAAM's eligibility ceiling is a program-design
threshold set by the nonprofit each year rather than a measurement of what its
members earn, and its member roster is a self-selected subset of the Austin
music workforce (those who applied for HAAM services), so the figure indexes
program eligibility rather than typical musician earnings. Some HAAM reports
also restate 2015 Austin Music Census figures alongside current impact data;
this study takes only the eligibility threshold from HAAM, not the reprinted
census numbers.} None of these three sources was designed as a representative
measurement of the Austin music workforce, and none is one. Each still adds
something the federal system does not collect: the purchased series reaches
self-employment earnings that payroll surveys exclude, the interview project
preserves how roughly sixty working musicians describe their finances in
their own words, and the HAAM ceiling marks where Austin's music-health
infrastructure draws its eligibility line. This study is built to complement
that Austin-specific work rather than replace it, by adding the policy-facing
records that update on a schedule: federal household microdata on earnings
and self-employment, the Comptroller's monthly venue receipts, and the city's
budget and awardee files.

\textbf{Austin's own music census has not asked a musician what they earn since
the 2014 fielding, whose income questions covered calendar 2013.} The city has
since relied on that 2013 figure, the three sources above, and federal series
never designed to capture irregular self-employed performance income.

\subsection{One quoted income figure is a sideman's 1981 salary}

When no current number exists, people quote whatever number they can find. The
clearest example is \$15,475, quoted as an Austin musician's annual income. It
originates in a February 2022 KUT public radio segment in which the bandleader
Marcia Ball described her own payroll from four decades earlier: ``if you had
worked for me for a full year as my sax player did, you would have made
\$15,475.''\footnote{Elizabeth McQueen and Miles Bloxson, ``Marcia Ball on the
economics of being a working musician in Austin then and now,'' KUT 90.5 and
KUTX, broadcast 15 February 2022. Marcia Ball is a Grammy-nominated Austin
bandleader whose regional touring band worked out of Austin in the early 1980s;
the \$15,475 figure is the annual salary she paid her saxophone player in 1981,
recalled from her own memory and payroll rather than measured by any survey.
The station's next sentence restated the figure in the present tense. A
full-text search of every document this study gathered found the number in
that KUT article and nowhere else; it appears in neither of the two Austin
music censuses and did not travel beyond its source. The KUT segment carried the number
forward to about \$47,865 in 2022 dollars; chaining the Bureau of Labor
Statistics Consumer Price Index for All Urban Consumers (series
CUUR0000SA0) from its 1981 annual average to its 2025 annual average puts the
1981 salary at roughly \$54,000 in 2025 dollars, which inverts the argument
the figure is usually used to support.}

A related claim about Austin journalism deserves the same scrutiny, even
though it favours this study's own case. Reporters and advocates often say
that Austin news coverage keeps recycling 2013 income data. This study
inventoried thirteen Austin news items published after the 2022 census that
touched on musician earnings or the music program, and that description is
too strong: only three of the thirteen cite an earnings figure at all, and
just one cites the 2013 census figure, under a subheading that flags the
figure's age.\footnote{This study
keeps that inventory in its Austin press-coverage evidence folder
(\texttt{01\_evidence/09\_austin\_press}), which the technical appendix
describes in full. It was assembled by keyword search across the archives of
KUT, KUTX, the \textit{Austin Chronicle}, the \textit{Austin American-Statesman},
\textit{Austin Monitor} and \textit{CultureMap Austin}, and it is a curated
sample of coverage rather than an exhaustive census. Any share drawn from it
is a share of the items this study gathered, not of Austin coverage
generally.} Coverage moved from earnings to program spending, reporting grant
totals, award sizes and city booking rates. The recycling did happen earlier:
two of the five 2022 items in the same inventory restated 2013 census income
statistics as current.

\subsection{Four federal series count this workforce, and each misses a different part of it}

\begin{sidebar}[title=What the federal statistics can and cannot see]
{\footnotesize
Picture a working Austin musician with a part-time day job, a dozen bar gigs a
month, a few students and some streaming royalties. Four federal statistical
programs each observe part of that person's working life, and none observes
all of it.

\begin{itemize}
\item \textbf{The wage survey asks employers, so it sees only payroll.} The
Bureau of Labor Statistics's Occupational Employment and Wage Statistics
program (OEWS) is a semi-annual establishment survey run jointly by BLS and
the state workforce agencies. Employers in the sample report the wage of every
worker on their payroll and report nothing for the self-employed, who fall
outside the sampling frame; the program describes its coverage at
\url{https://www.bls.gov/oes/oes_ques.htm}. OEWS records
\NEmpPayrollMusiciansAustinLast\ wage-and-salary musician jobs in the Austin
metro in 2025, against an estimated \NPumsAustinMusWtdPos\ people with positive
earnings whom the American Community Survey records as working as musicians
here over 2020 to 2024. Across every musician row in the OEWS extract used
here, covering Austin, Texas and the nation from 2005 to 2025, the program
publishes an hourly wage and no annual one.\footnote{Annual mean and annual
median are blank in all 63 musician rows of that extract, in every year and
every geography. The Bureau states this as a standing convention for
occupations with substantial part-year work rather than as intermittent
suppression; the wage series therefore cannot be annualised from within its
own tables. Any annual figure for Austin musicians built from OEWS would
require assumptions about weeks and hours worked that OEWS itself does not
supply, and this study does not construct one.}

\item \textbf{The quarterly payroll census is built from unemployment-insurance
records.} The BLS Quarterly Census of Employment and Wages (QCEW) tabulates
every job covered by state unemployment-insurance law, drawing directly from
the wage records employers file with the state; program page at
\url{https://www.bls.gov/cew/}. Because it is UI-covered employment only, an
unincorporated self-employed performer is outside QCEW coverage by
construction. Performers who incorporate as an S-corp or LLC and pay themselves
wages do enter QCEW, which is part of why the independent-artists industry code
(NAICS 7115) shows any employment at all.

\item \textbf{Nonemployer Statistics does see the self-employed}, and counts
their businesses. The Census Bureau builds Nonemployer Statistics (NES) from
federal tax returns filed by businesses with no paid employees and receipts of
at least \$1,000, and publishes county-by-industry counts and gross receipts
at \url{https://www.census.gov/programs-surveys/nonemp.html}. It reports
receipts before any expense, and no federal source records what a performer
spends to earn them. Separately, the 501 respondents to the 2022 Greater
Austin Music Census who answered every spending category report a median of
\NCensusSpendTotalMedian\ a year spent on their own music
work.\footnote{Median from this study's own tabulation of the census
microdata, dataset \texttt{an3p-3yqx} at
\url{https://data.austintexas.gov/resource/an3p-3yqx.csv}, deflated to 2025
dollars with the CPI-U. The 2022 census is a self-selected online convenience
sample with no survey weights, so this median describes the 501 respondents
who answered the item rather than any wider population; it counts people
rather than businesses, so it cannot be netted against NES receipts.}

\item \textbf{Only the American Community Survey asks people directly} and
adds self-employment income to wages. The Census Bureau's American Community
Survey (ACS) is a continuous household survey with about 3.5 million
respondents a year nationwide; the Bureau releases a Public Use Microdata
Sample (PUMS) at
\url{https://www2.census.gov/programs-surveys/acs/data/pums/2024/5-Year/} that
lets analysts identify a respondent's detailed occupation, add wage and
self-employment earnings on the same record, and attach person weights and 80
replicate weights for sampling error. Even ACS records one occupation per
person, the current or most recent job, so a musician with a day job is
counted as whatever the day job is.
\end{itemize}

\smallskip
The federal system measures this workforce continuously and undercounts it
systematically. The more of a musician's income comes from gigs, teaching and
royalties rather than from a payroll job, the less any of these series describes
that musician.}
\end{sidebar}

The pattern already shows up in Austin's own numbers. Two hourly-wage series
for the same occupation, in the same metro, in the same year disagree by
almost half. The City of Austin's purchased Creative Vitality Suite figure for
Musicians and Singers in the Austin metro in 2017 was \$17.68 an hour, against
\$31.03 an hour in the BLS Occupational Employment and Wage Statistics
tabulation for the same occupation, metro and year: the purchased figure came
in at \NColAnchorCityWageVsOewsTwoZeroOneSeven\ of the federal
one.\footnote{City figure from dataset \texttt{2qxc-8cme} at
\url{https://data.austintexas.gov/resource/2qxc-8cme.csv}; federal figure from
the Occupational Employment and Wage Statistics May 2017 estimates at
\url{https://www.bls.gov/oes/tables.htm}, SOC 27-2042, Austin-Round
Rock-San Marcos MSA. Every dollar in this comparison is a same-year nominal
figure, so no deflation applies.} The federal series counts payroll jobs, and
the purchased series appears to combine payroll wages with self-employment
earnings, though the vendor publishes no methodology, so the size of the gap
is a weaker piece of evidence than its direction. Section~\ref{sec:costs}
returns to that disagreement when it asks what a one-bedroom apartment costs
an Austin musician.

The \NEmpPayrollMusiciansAustinLast\ Austin jobs behind the federal series in
2025 are also a small base for a wage estimate, and that base size matters
wherever this study uses the payroll wage: the Austin musician wage moves by
several dollars an hour between adjacent years, so its level in any single
year is uncertain. This study still uses the OEWS series, because nothing else
records what payroll music jobs pay in Austin across two decades, and we
suggest reading its direction over a run of years rather than any single
year's value.

\subsection{How this study keeps mismatched sources apart}

The rules below apply to every number that follows: every figure carries a
named source, the population behind it is stated, and sources that measure
different things stay in separate sentences.

\begin{itemize}
\item Every dollar figure is adjusted for inflation to 2025 dollars using the
      Bureau of Labor Statistics Consumer Price Index for All Urban Consumers,
      U.S.\ city average, all items (series CUUR0000SA0), and the base year
      appears in every figure subtitle. Three kinds of figure are left nominal
      and are flagged as such where they appear: a ratio of two same-year
      dollar figures, where deflation cancels; a figure quoted as its source
      originally published it; and award totals whose payments all fall inside
      one program period.\footnote{BLS publishes a bimonthly CPI for
      Dallas-Fort Worth-Arlington, a bimonthly CPI for Houston-The Woodlands-Sugar
      Land and a monthly South regional series, but no separate CPI for Austin,
      so this study deflates Austin dollar figures with the national CUUR0000SA0
      series retrieved via the BLS Public Data API at
      \url{https://api.bls.gov/publicAPI/v2/timeseries/data/}. No measurement
      in this study establishes whether Austin prices rose faster or slower
      than the national index over the period, so the direction of that
      substitution's effect on the estimates is unknown. October 2025 CPI
      collection was suspended during the federal funding lapse; the missing
      month is linearly interpolated in this project's file and flagged in the
      figures that touch it.}
\item Every American Community Survey estimate is reported with a margin of
      error computed by the Census Bureau's successive-difference replication
      method, using the 80 alternative person weightings the Bureau publishes
      alongside the main person weight (variables PWGTP1 through PWGTP80) so
      users can measure sampling error; the method and the replicate-weight
      file are documented in the Bureau's PUMS Accuracy of the Data at
      \url{https://www2.census.gov/programs-surveys/acs/tech_docs/pums/accuracy/}.
      This study publishes no ACS estimate drawn from fewer than 50 unweighted
      records.
\item The 2022 Greater Austin Music Census is a self-selected online survey
      with no survey weights, so its percentages describe the 2,227 people who
      answered it rather than the Austin music workforce as a whole; no margin
      of error is computable for any of them, and a probability-sample
      benchmark for the same population does not exist. Its shares never appear
      in the same comparison as an American Community Survey estimate.
\end{itemize}

Those rules cost the study some tidy comparisons. They are also why every
Austin musician estimate in the next section appears with its margin of error
attached.
